%======================================================================
\chapter{Non-cut points}\label{ch:noncut}
%======================================================================

In this chapter $X$ is a topological space, not necessarily metric. The main theorem is due to
R.~L.~Moore; we follow Nadler (\emph{Continuum Theory}, 6.3--6.7), as the formal proof does. The
space will be compact, connected and $T_1$ (points are closed). It is applied later not to a
continuum but to the decomposition space of a decomposition of a continuum
(Remark~\ref{rem:noncut-use}), which is why we do not assume a metric.

\section{Separations and cut points}

\begin{definition}\label{def:separation}
Let $Y\subseteq X$. A \emph{separation} $Y=P\mid Q$ is a pair of nonempty disjoint sets $P,Q$ with
$P\cup Q=Y$, both open in the subspace $Y$. A point $x\in X$ is a \emph{cut point} of $X$ if
$X\setminus\set x$ is not connected; otherwise it is a \emph{non-cut point}.
\end{definition}

\formal{\lean{HPC.Plane.IsSeparation}, \lean{HPC.Plane.IsCutPoint} (in
\lean{HPC/Plane/NonCut51Gen.lean}).}

\begin{lemma}\label{lem:sep-basic}
Let $Y=P\mid Q$ be a separation.
\begin{enumerate}[label=(\alph*)]
\item $P$ and $Q$ are separated: $\cl P\cap Q=\emptyset=P\cap\cl Q$.
\item If $Y$ is open in $X$, then $P$ and $Q$ are open in $X$.
\item A preconnected subset of $Y$ lies in $P$ or in $Q$.
\item If $Y$ is nonempty and not connected and $c\in Y$, then $Y$ has a separation $Y=P\mid Q$ with
$c\in Q$.
\end{enumerate}
\end{lemma}

\begin{proof}
(a) Write $Q=Y\cap t$ with $t$ open in $X$. Then $t\cap P=\emptyset$, so $t\cap\cl P=\emptyset$, so
$Q\cap\cl P=\emptyset$. Symmetrically for the other half.
(b) $Q=Y\cap t$ is an intersection of two open sets.
(c) is Lemma~\ref{lem:side} (whose proof uses only the topology), with (a).
(d) As $Y$ is not preconnected there are open $u,v\subseteq X$ with $Y\subseteq u\cup v$, $Y\cap u\cap
v=\emptyset$ and both $Y\cap u$, $Y\cap v$ nonempty. These form a separation; name them so that
$c$ lies in the second.
\end{proof}

\formal{\lean{HPC.Plane.IsSeparation.disjoint_closure_left},
\lean{HPC.Plane.IsSeparation.isOpen_right}, \lean{HPC.Plane.IsSeparation.subset_or_subset},
\lean{HPC.Plane.exists_separation_of_not_isConnected}.}

\begin{theorem}\label{thm:nadler63}
Let $X$ be connected, $C\subseteq X$ connected, and $X\setminus C=A\mid B$. Then $A\cup C$ is
connected.
\end{theorem}

\begin{proof}
$A\cup C$ is nonempty. Let $u,v$ be open sets with $A\cup C\subseteq u\cup v$ and $(A\cup C)\cap u\cap
v=\emptyset$; we must show $A\cup C\subseteq u$ or $A\cup C\subseteq v$. Since $C$ is connected, it lies
in one of them, say $C\subseteq u$. Put
\[
P=A\cap v,\qquad W=(A\cap u)\cup C .
\]
Then $W\subseteq u$, and $X=P\cup W\cup B$ (a point not in $C$ lies in $A$ or $B$, and a point of $A$ lies
in $u$ or $v$). The sets $P$ and $W\cup B$ are separated:
\begin{itemize}
\item $\cl P\cap W=\emptyset$: the open set $u$ misses $P$ (a point of $P\cap u$ would be in $(A\cup
C)\cap u\cap v$), hence misses $\cl P$, and $W\subseteq u$.
\item $\cl P\cap B\subseteq\cl A\cap B=\emptyset$ by Lemma~\ref{lem:sep-basic}(a).
\item $P\cap\cl W=\emptyset$: the open set $v$ misses $W$ (as $W\subseteq(A\cup C)\cap u$), hence misses
$\cl W$, and $P\subseteq v$.
\item $P\cap\cl B\subseteq A\cap\cl B=\emptyset$.
\end{itemize}
$B$ is nonempty and $X$ is connected, so $P=\emptyset$ by Lemma~\ref{lem:side}. Hence $A\subseteq u$ and
$A\cup C\subseteq u$.
\end{proof}

\formal{\lean{HPC.Plane.isConnected_union_of_separation} (Nadler 6.3).}

\begin{lemma}\label{lem:nadler64}
Let $X$ be connected, $X\setminus\set x=K\mid L$ and $X\setminus\set y=M\mid N$, with $x\in M$ and
$y\in K$. Then $N\cup\set y\subseteq K$.
\end{lemma}

\begin{proof}
$x\ne y$ since $x\in M\subseteq X\setminus\set y$. By Theorem~\ref{thm:nadler63} (with $C=\set y$ and the
separation $X\setminus\set y=N\mid M$) the set $N\cup\set y$ is connected. It lies in $X\setminus\set
x$: $x\in M$ is not in $N$, and $x\neq y$. By Lemma~\ref{lem:sep-basic}(c) it lies in $K$ or in $L$; it
contains $y\in K$, so it lies in $K$.
\end{proof}

\formal{\lean{HPC.Plane.nadler_6_4}.}

\section{Moore's theorem}

\begin{theorem}[Moore]\label{thm:noncut-side}
Let $X$ be a compact connected $T_1$ space, $c\in X$, and $X\setminus\set c=U\mid V$. Then $U$
contains a non-cut point of $X$, and so does $V$.
\end{theorem}

\begin{proof}
By symmetry it suffices to treat $U$. Suppose every point of $U$ is a cut point. For each $x\in U$,
Lemma~\ref{lem:sep-basic}(d) gives a separation $X\setminus\set x=U_x\mid V_x$ with $c\in V_x$ (note
$c\neq x$). Put $A(x)=U_x\cup\set x$. We prove four facts.

\emph{(i) $A(x)\subseteq U$.} Apply Lemma~\ref{lem:nadler64} to the separations $X\setminus\set
c=U\mid V$ and $X\setminus\set x=V_x\mid U_x$: we have $c\in V_x$ and $x\in U$, so $U_x\cup\set x\subseteq
U$.

\emph{(ii) $A(x)$ is closed.} Its complement is $V_x$, which is open in $X$ by
Lemma~\ref{lem:sep-basic}(b), since $X\setminus\set x$ is open ($X$ is $T_1$).

\emph{(iii) If $y\in U_x$, then $A(y)\subseteq U_x$; in particular $x\notin A(y)$.} By (i), $y\in U$. Since
$x\neq y$, $x$ lies in $U_y$ or in $V_y$. If $x\in U_y$, Lemma~\ref{lem:nadler64} applied to
$X\setminus\set y=U_y\mid V_y$ and $X\setminus\set x=U_x\mid V_x$ (with $y\in U_x$ and $x\in U_y$) gives
$V_x\cup\set x\subseteq U_y$; but $c\in V_x$ and $c\in V_y$, a contradiction. So $x\in V_y$, and
Lemma~\ref{lem:nadler64} applied to $X\setminus\set x=U_x\mid V_x$ and $X\setminus\set y=V_y\mid U_y$
(with $x\in V_y$ and $y\in U_x$) gives $U_y\cup\set y\subseteq U_x$. Finally $x\notin U_x$.

\emph{(iv) The family $\set{A(x):x\in U}$ has a minimal member.} We use Zorn's lemma in the form stated before
Lemma~\ref{lem:nested-meets}. The family is nonempty ($U\neq\emptyset$). Let $\mathcal C$ be a
nonempty chain in it. Its members are closed and nonempty, any finitely many have nonempty
intersection, and $X$ is compact; so there is a point $p$ in all of them. Then $p\in U$ by (i). For
each member $A(x)$ of $\mathcal C$, either $p\in U_x$, and then $A(p)\subseteq U_x\subseteq A(x)$ by (iii),
or $p=x$ and $A(p)=A(x)$. So $A(p)$ is a member of the family contained in $\bigcap\mathcal C$.

Now let $A(x)$ be minimal, and take $q\in U_x$ (it is nonempty). By (iii) $A(q)\subseteq U_x\subseteq
A(x)$, so $A(q)=A(x)$ by minimality, and $x\in A(x)=A(q)$, contradicting (iii).
\end{proof}

\formal{\lean{HPC.Plane.exists_nonCutPoint_left} and \lean{HPC.Plane.exists_nonCutPoint_of_cut}
(Nadler 6.6). Nadler assumes $X$ nondegenerate; this follows from the separation and is not
needed.}

\begin{theorem}\label{thm:two-noncut}
A compact connected $T_1$ space with at least two points has two distinct non-cut points.
\end{theorem}

\begin{proof}
If $X$ has no cut point, any two distinct points will do. If $c$ is a cut point, take $d\ne c$ and a
separation $X\setminus\set c=U\mid V$ (Lemma~\ref{lem:sep-basic}(d)); Theorem~\ref{thm:noncut-side}
gives non-cut points $u\in U$ and $v\in V$, which are distinct as $U\cap V=\emptyset$.
\end{proof}

\formal{\lean{HPC.Plane.exists_two_nonCutPoints}.}

\begin{theorem}\label{thm:noncut-not-in}
Let $X$ be a compact connected $T_1$ space and $Z\subsetneq X$ a connected proper subset. Then some
non-cut point of $X$ lies outside $Z$.
\end{theorem}

\begin{proof}
Suppose all non-cut points lie in $Z$. Take $c\notin Z$; it is a cut point. Take $z\in Z$ and a
separation $X\setminus\set c=U\mid V$. The connected set $Z\subseteq X\setminus\set c$ lies in $U$ or in
$V$, say in $U$. By Theorem~\ref{thm:noncut-side} $V$ contains a non-cut point, which is then in
$Z\subseteq U$; but $U\cap V=\emptyset$.
\end{proof}

\formal{\lean{HPC.Plane.not_subset_of_nonCutPoints} (Nadler 6.7).}

\begin{example}
The non-cut points of $[0,1]$ are $0$ and $1$. Every point of a circle is a non-cut point. In the
letter T (three segments with a common end point $o$), the non-cut points are the three free end
points; $o$ is a cut point whose complement has three components.
\end{example}

\begin{remark}\label{rem:noncut-use}
In the later volumes (Bing's theorem on atriodic hereditarily decomposable continua) these
results are applied to the decomposition space of an upper semicontinuous decomposition of a
continuum into subcontinua (Chapter~\ref{ch:decomp}): the decomposition space is compact,
connected and $T_1$, and its non-cut points correspond to the ``end elements'' of the
decomposition (\lean{HPC.Plane.exists_two_endElts}).
\end{remark}

\section*{Exercises}

\begin{exercise}
Show that compactness cannot be dropped from Theorem~\ref{thm:two-noncut}: every point of $\R$ is
a cut point. Where does the proof of Theorem~\ref{thm:noncut-side} use compactness, and where
does it use the $T_1$ property?
\end{exercise}

\begin{exercise}
Show that a nondegenerate continuum with exactly two non-cut points is an arc. (Hard; this is a
classical characterization of the arc. It is not used in these notes.)
\end{exercise}

\begin{exercise}
Let $X$ be a compact connected $T_1$ space and $c$ a cut point. Show that each component of
$X\setminus\set c$ contains a non-cut point of $X$, if $X\setminus\set c$ has finitely many components.
\end{exercise}
